\section{Proof of the transfer theorem}
\begin{lemma}[Conditional task comparison]\label{lem:comparison}
For every legal pre-audit history, report and prediction $r\in H$,
\[
 k_-\{w_a(y)+\|f_a(y)-r\|^2\}
 \le \E[\|Z-r\|^2\mid\mathcal H_{t-1},a,y]-B_0
 \le k_+\{w_a(y)+\|f_a(y)-r\|^2\}.
\]
The same inequalities remain valid after an independent routing coin chooses $r$.
\end{lemma}
\begin{proof}
Write $m=m_{t,a}(y)$, $v=v_{t,a}(y)$ and $w=w_a(y)$. Conditional projection gives the middle expression as $v-B_0+\|m-r\|^2$. The upper follows from
$\|m-r\|^2\le2\|f_a-r\|^2+2c_bw$.
For the lower,
\[
 w+\|f_a-r\|^2\le(1+2c_b)w+2\|m-r\|^2
 \le k_-^{-1}\{v-B_0+\|m-r\|^2\}.
\]
The coin has no audit information, so conditioning on its selected value does not change $m$ or $v$.
\end{proof}

\begin{lemma}[Actual stopped occupation]\label{lem:occupation}
For any legal history-dependent procedure of finite expected cost, put
\[
 x_a=\E\sum_{t\le\tau}\one_{\{A_t=a\}},\qquad
 \zeta_a(D)=\Pp(A_\tau=a,Y_\tau\in D).
\]
Then $\sum_a\zeta_a(Y_a)=1$, $\sum_ac_ax_a\le N$, and
$\zeta_a\le u_ax_a\nu_a$. These are assertions about its actual execution law.
\end{lemma}
\begin{proof}
Positive minimum cost implies $\E\tau<\infty$ and hence almost-sure termination. For a Borel set $D$, the expected number of action-$a$ reports in $D$ is
\[
 \sum_t\E\left[\one_{\{\tau\ge t,A_t=a\}}
       \Pp(Y_t\in D\mid\mathcal H_{t-1},A_t=a)\right]
 \le u_ax_a\nu_a(D).
\]
The terminal count is no larger than this acquired count. Predictability is used before the report; conditioning the law after rejection back to a reference law would not justify this calculation.
\end{proof}

\begin{proof}[Lower bounds]
Apply Lemma~\ref{lem:comparison} to the actually committed output, then replace its squared distance by the minimum over the fixed output dictionary. With $h_a=d\zeta_a/d\nu_a$ and $x'_a=u_ax_a$, Lemma~\ref{lem:occupation} yields a feasible point of $\Xi_m(N;c/u)$. This proves the checkpoint lower for each model and then for the common worst-model infimum. An autonomous program has at least one live label, so at most $M-1$ terminal values. Its entire action and stopping policy is already included in the preceding lower bound.
\end{proof}

\begin{proof}[Common finite-state upper]
Take a minimizing point $(C,x,h)$ of $\Xi_m(N;c/\ell)$ and put $T=\sum_a x_a/\ell_a\ge1$. At one live label, choose action $a$ with probability $x_a/(\ell_aT)$. For $x_a>0$, accept report $y$ with probability $h_a(y)/x_a$, route to a nearest proxy center in $C$, and otherwise continue at the same live label. Actions with $x_a=0$ are not selected. A fixed Borel tie rule is used.

At every surviving history, the conditional probability of stopping is at least
$T^{-1}\sum_a\int h_a\,d\nu_a=T^{-1}$. Therefore $\E_\theta\tau\le T$, uniformly in $\theta$. Conditional expected acquisition cost per call is the fixed number $T^{-1}\sum_ac_ax_a/\ell_a$, so total expected cost is at most $N$. The conditional contribution to protected excess risk on one call is at most
\[
 \frac{k_+}{T}\sum_a\frac{u_a}{\ell_a}\int d_{a,C}h_a\,d\nu_a
 \le \frac{k_+K_0}{T}\Xi_m(N;c/\ell).
\]
Summing over actual survival probabilities proves the upper. It uses no persistent real vector, posterior or random seed. The parameter-independent program has one live label and $m$ terminal labels; its actual selected law need not equal $h_a\nu_a$.
\end{proof}

\begin{proof}[Existence, dual and exact special case]
The counts belong to $0\le x_a\le N/c_a$. On this box the functions $h_a$ lie in weak-star compact subsets of $L^\infty(\nu_a)$. The inequalities $0\le h_a\le x_a$ and the mass constraint are weak-star closed, as can be tested against nonnegative $L^1$ functions. The centers belong to compact $H^m$, and their squared distance functions converge uniformly when the centers converge. Since $w_a$ is bounded, the objective is jointly continuous in these topologies. This proves existence and the feasibility criterion.

For fixed centers, minimizing the mass--cost Lagrangian over $0\le h_a\le x_a$ gives
\[
 z-\lambda N+\sum_ax_a\left\{\lambda c_a-\int(z-d_{a,C})_+d\nu_a\right\}.
\]
Minimizing in $x_a\ge0$ yields \eqref{eq:dual}. To prove equality, consider the convex cone of achievable triples (terminal mass, cost, marked distortion), with nonnegative slack in the last two coordinates. Its bounded-cost slices are closed by the compactness just proved. At mass one and $N>\min c_a$, the projected point is interior to the feasible mass--cost cone. A supporting plane to the finite convex distortion value has a nonpositive cost coefficient and can be normalized to distortion coefficient one. Testing it on the one-action rays gives precisely \eqref{eq:dual}, with equality at the optimum. At the endpoint, take budgets decreasing to $\min c_a$; compactness gives convergence of the primal values, and evaluation of the interior dual optimizers at the endpoint gives the supremum formula. No boundedness of endpoint multipliers is inferred.

For the two-action assertion, at an optimizer discard actions with no terminal mass. Freeze each remaining action's conditional terminal distribution, per-unit terminal cost and per-unit terminal distortion. Reoptimize only its weights subject to total mass one and one cost inequality. An extreme minimizer has at most two positive weights: three positive weights admit a nonzero two-sided perturbation preserving total mass and cost. Rescaling the frozen counts and densities preserves their pointwise domination. This standard finite-dimensional argument makes no claim about a uniquely optimal action mixture.

Finally, in the zero-mark regenerative case the conditional loss is exactly $B_0+\|f_a-r\|^2$. The lower occupation inequality uses $u_a=1$. The construction chooses $a$ with probability $x_a/\sum x_a$ and accepts with $h_a/x_a$; its stopping mass is exactly $1/\sum x_a$ at every call. Its actual terminal submeasures are exactly $h_a\nu_a$, and its expected action counts are exactly $x_a$. The exact equalities follow. This last mechanism is the inherited report-recoverable result~\cite{GTF51}; the conditional task comparison is what permits the nonzero-mark extension.
\end{proof}
